% se puede agregar la opción [english] para 
%  memorias o tesis en inglés (borrando el archivo .aux)
\documentclass[draft]{umemoria} 

\usepackage[color]{theoremenvFMG}
\usepackage[swapeps,swapphi]{mathcmdFMG}

\depto{Departamento de ???}
\author{Nombre Completo Autor}
\title{Título de la Memoria/Tesis}

% incluir ambos comandos para una doble titulación
%  o quitar el comando que no aplica
\memoria{Ingenier[o/a?] Civil en ???}
\tesis{Magíster en ???}
%\tesis{Doctor en ???} % incluir solo este comando para doctorados

% puede haber varios profesores guía separados por coma;
% pero si es una memoria, solo puede haber un profesor guía
\guia{Nombre Completo Guía} 

% puede haber varios profesores co-guía separados por coma;
% pero si es una memoria, el profesor co-guía será el primer
% integrante de la comisión
%\coguia{Nombre Completo Co-Guía} % incluir en caso de co-guía de *tesis*

%\cotutela{Nombre Institución} % incluir en caso de cotutela

\comision{Nombre Completo Uno,Nombre Completo Dos,Nombre Completo Tres}

%\auspicio{Nombre Institución} % incluir en caso de recibir financiamiento

% tiene que ser el año en que se da el examen de título/grado (defensa)
%\anho{2021} % incluir solo para reemplazar el año actual

\usepackage{lipsum}

% \includeonly{lib/cap2.tex,lib/cap3.tex}

\begin{document}

\frontmatter
\maketitle

\begin{resumen}
    \lipsum[1-4]
\end{resumen}

% opcional: incluir para tesis en inglés;
%  en este caso hay que tener el resumen y abstract
%   en ambos idiomas
%\begin{abstract}
%\lipsum[1-4]
%\end{abstract}

\begin{dedicatoria}
    Una dedicatoria corta.
\end{dedicatoria}

\begin{thanks}
    \lipsum[1-2]
\end{thanks}

\tableofcontents
\listoftables % opcional
\listoffigures % opcional

\mainmatter

\chapter{Introducción}

\lipsum[30-35]
\begin{enumerate}
    \item Item 1
          \begin{enumerate}
              \item Subitem 1
              \item Subitem 2 (ver Figura \ref{logofcfm})
          \end{enumerate}
    \item Item 2
    \item Item 3
\end{enumerate}
\begin{theorem}
    Se tiene que $$\int_0^t e^s \dd{s}=e^t-1.$$
    $\R, \N, \Q$,
    $\epsilon, \varepsilon$
    $\phi, \varphi$
\end{theorem}
\begin{proof}
    Demostración asombrosa.
\end{proof}
\lipsum[36-40]
\chapter{Segundo}

\lipsum[1-3]

\begin{definition}[ver \cite{KAR00}]
    Definición definitiva $$\frac{d}{dx}\int_a^xf(y)dy=f(x).$$
\end{definition}
\chapter{Tercero}
\lipsum[50-60]
\chapter{Conclusión}

\lipsum[130-132]
\begin{figure}
    \centering
    \includegraphics[scale=.2]{img/logos/fcfm.pdf}
    \caption{Logo de la Facultad}
    \label{logofcfm}
\end{figure}

\lipsum[133-134]

\begin{table}
    \centering
    \caption{Tabla 1}
    \label{tabla:1}
    \begin{tabular}{|c|c|r|}
        \hline
        \textbf{Campo 1} & \textbf{Campo 2} & \textbf{Num} \\\hline
        Valor 1a         & Valor 2a         & 3            \\
        Valor 1b         & Valor 2b         & 3            \\
        \hline
    \end{tabular}

\end{table}

\lipsum[135]


% ver https://www.overleaf.com/learn/latex/Glossaries
% \input{lib/glosario.tex} % opcional

\nocite{*}
\bibliographystyle{plain}
\bibliography{bibliografia}

% opcional ...
\begin{appendices}
    \chapter{Anexo}
\lipsum[50-60]
\end{appendices}
\end{document}